% !TEX program = xelatex
\documentclass[10pt,xcolor=dvipsnames]{beamer}
\usepackage{fontspec}
\usepackage{xeCJK}
\usepackage{polyglossia}
\usepackage{url}
\setdefaultlanguage{vietnamese}

% Fonts
\setmainfont{Linux Libertine O}[
  BoldFont=Linux Libertine O Bold,
  ItalicFont=Linux Libertine O Italic
]
\setsansfont{Linux Biolinum O}[
  BoldFont=Linux Biolinum O Bold
]
% \setmonofont{Inconsolata}  % Font không có sẵn, dùng fallback

% Theme
\usetheme{metropolis}
\usecolortheme{default}

% Colors
\definecolor{HUEBlue}{RGB}{0,51,102}
\definecolor{HUEGold}{RGB}{212,175,55}
\setbeamercolor{frametitle}{bg=HUEBlue,fg=white}
\setbeamercolor{title}{fg=HUEBlue}
\setbeamercolor{structure}{fg=HUEBlue}

% Remove navigation symbols
\setbeamertemplate{navigation symbols}{}

\begin{document}

% Title page - custom design
\begin{frame}
  \begin{center}
    % Logos
    \includegraphics[height=1.2cm]{logo_thpt_thuan_hoa.png}
    \hfill
    \includegraphics[height=1.2cm]{logo_dhsp_hue.png}
    
    \vspace{0.2cm}
    {\scriptsize \textbf{TRƯỜNG ĐẠI HỌC SƯ PHẠM, ĐẠI HỌC HUẾ \quad -- \quad TRƯỜNG THPT THUẬN HOÁ}}
    
    \vspace{0.3cm}
    {\color{HUEBlue} \rule{0.9\textwidth}{0.4pt}}
    
    \vspace{0.2cm}
    {\large \textbf{\textcolor{HUEBlue}{THUYẾT MINH ĐỀ CƯƠNG NGHIÊN CỨU}}}
    
    \vspace{0.2cm}
    {\small \textbf{XÂY DỰNG PHÒNG THÍ NGHIỆM AI OFFLINE (AINATOMY)}\\[0.1em]
    \textit{Phát triển năng lực kiểm chứng AI cho học sinh THPT theo CT GDPT 2018}}
    
    \vspace{0.2cm}
    {\color{HUEBlue} \rule{0.9\textwidth}{0.4pt}}
    
    \vspace{0.2cm}
    {\small
    \begin{tabular}{ll}
      \textbf{Chủ trì:} & Mai Xuân Văn -- GV Tin học \\
      \textbf{Đơn vị:} & Trường THPT Thuận Hoá \\
      \textbf{Phối hợp:} & Trường Đại học Sư Phạm, Đại học Huế \\
    \end{tabular}
    }
    
    \vspace{0.2cm}
    {\scriptsize \textit{Giải thưởng Tiên phong Ứng dụng AI trong Giáo dục Việt Nam 2026 \quad -- \quad Năm học 2025--2026}}
  \end{center}
\end{frame}

% Slide 2: Outline
\begin{frame}{Mục lục}
  \tableofcontents
\end{frame}

% Section 1: Mở đầu
\section{Mở đầu}

\begin{frame}{Lý do chọn đề tài}
  \begin{block}{Chính sách giáo dục}
    Quyết định 2422/QĐ-BGDĐT yêu cầu 12 tiết/năm cho giáo dục AI
  \end{block}
  
  \begin{block}{Thách thức hạ tầng}
    Phòng máy yếu, internet không ổn định ở vùng khó khăn
  \end{block}
  
  \begin{block}{Khoảng trống năng lực}
    Học sinh thiếu kỹ năng kiểm chứng đầu ra AI
  \end{block}
\end{frame}

\begin{frame}{Mục tiêu nghiên cứu}
  \textbf{Mục tiêu tổng quát:} Xây dựng AInatomy giúp học sinh phát triển năng lực kiểm chứng AI, đáp ứng CT GDPT 2018 và QĐ 2422.
  
  \vspace{1em}
  \textbf{Mục tiêu cụ thể:}
  \begin{enumerate}
    \item Thiết kế công cụ offline, không cần tài khoản
    \item Xây dựng 28 bài học theo 22 YCCĐ Tin học 10
    \item Đánh giá hiệu quả phát triển năng lực kiểm chứng AI
  \end{enumerate}
\end{frame}

\begin{frame}{Đối tượng \& Phạm vi}
  \begin{columns}[T]
    \begin{column}{0.48\textwidth}
      \textbf{Đối tượng:}
      \begin{itemize}
        \item Công cụ AInatomy
        \item Học sinh THPT (lớp 10)
        \item Giáo viên Tin học
        \item CT GDPT 2018, Tin học 10
      \end{itemize}
    \end{column}
    
    \begin{column}{0.48\textwidth}
      \textbf{Phạm vi:}
      \begin{itemize}
        \item 28 bài học, 22 YCCĐ
        \item Trường THPT Thuận Hoá
        \item Năm học 2025--2026
        \item HTML5, CSS3, JS, Three.js
      \end{itemize}
    \end{column}
  \end{columns}
\end{frame}

\begin{frame}{Phương pháp nghiên cứu}
  \begin{enumerate}
    \item \textbf{Nghiên cứu tài liệu:} Phân tích CT GDPT 2018, QĐ 2422
    \item \textbf{Thiết kế \& phát triển:} Quy trình lặp, kiểm thử từng tính năng
    \item \textbf{Thực nghiệm sư phạm:} Triển khai trong lớp học Tin học 10
    \item \textbf{Đánh giá \& phân tích:} Định lượng + định tính
  \end{enumerate}
\end{frame}

% Section 2: Nội dung
\section{Nội dung}

\begin{frame}{Bối cảnh \& Vấn đề}
  \begin{alertblock}{Khó nhìn thấy AI}
    AI là hộp đen, học sinh chỉ thấy đầu ra
  \end{alertblock}
  
  \begin{alertblock}{Khó triển khai}
    Phòng máy yếu, không thể yêu cầu tài khoản/API
  \end{alertblock}
  
  \begin{alertblock}{Khó kiểm chứng}
    Học sinh tin tuyệt đối vào AI, không hỏi lại
  \end{alertblock}
\end{frame}

\begin{frame}{Giải pháp: AInatomy}
  \textbf{Phòng thí nghiệm AI offline} giúp học sinh quan sát quy trình AI từ đầu đến cuối.
  
  \vspace{1em}
  \textbf{4 bước học tập:}
  \begin{enumerate}
    \item \textbf{Đọc:} Nhận diện tình huống
    \item \textbf{Dự đoán:} Dự đoán đầu ra AI
    \item \textbf{Thực hiện:} Chạy mô phỏng
    \item \textbf{Phản hồi:} So sánh và rút bài học
  \end{enumerate}
\end{frame}

\begin{frame}{Tính năng cốt lõi}
  \begin{columns}[T]
    \begin{column}{0.32\textwidth}
      \textbf{Nhà máy AI}
      \begin{itemize}
        \item Thu thập dữ liệu
        \item Huấn luyện
        \item Triển khai
        \item Đánh giá
      \end{itemize}
    \end{column}
    
    \begin{column}{0.32\textwidth}
      \textbf{Đấu trường bắt lỗi}
      \begin{itemize}
        \item Nhận đầu ra AI
        \item Kiểm tra lỗi
        \item Chấm điểm
      \end{itemize}
    \end{column}
    
    \begin{column}{0.32\textwidth}
      \textbf{Bản đồ năng lực}
      \begin{itemize}
        \item Theo dõi tiến độ
        \item Điểm mạnh/yếu
        \item 5 năng lực AI
      \end{itemize}
    \end{column}
  \end{columns}
\end{frame}

\begin{frame}{Kết quả \& Tác động}
  \begin{block}{Chỉ số đạt được}
    \begin{itemize}
      \item \textbf{28/28} YCCĐ chương trình Tin học 10
      \item \textbf{97/97} kiểm thử tự động đạt
      \item \textbf{100\%} offline, không cần internet
      \item \textbf{3D} mô phỏng bằng Three.js
    \end{itemize}
  \end{block}
  
  \vspace{0.5em}
  \textbf{Tác động:} Học sinh phát triển năng lực kiểm chứng AI, hiểu quy trình, sử dụng AI có trách nhiệm.
\end{frame}

\begin{frame}{Tính bền vững \& Nhân rộng}
  \begin{columns}[T]
    \begin{column}{0.32\textwidth}
      \textbf{Bền vững}
      \begin{itemize}
        \item Offline, không phụ thuộc internet
        \item Tương thích máy cấu hình thấp
        \item Mã nguồn mở
      \end{itemize}
    \end{column}
    
    \begin{column}{0.32\textwidth}
      \textbf{Nhân rộng}
      \begin{itemize}
        \item Mọi trường THPT triển khai được
        \item Không cần đào tạo phức tạp
        \item Tích hợp bài giảng hiện có
      \end{itemize}
    \end{column}
    
    \begin{column}{0.32\textwidth}
      \textbf{Phát triển}
      \begin{itemize}
        \item Thêm bài học nâng cao
        \item Tích hợp AI thực tế
        \item Nghiên cứu tác động dài hạn
      \end{itemize}
    \end{column}
  \end{columns}
\end{frame}

% Section 3: Kết luận
\section{Kết luận}

\begin{frame}{Kết luận}
  \begin{itemize}
    \item AInatomy đáp ứng 28/28 YCCĐ CT GDPT 2018
    \item Chạy offline, phù hợp hạ tầng trường THPT Việt Nam
    \item Giúp học sinh phát triển năng lực kiểm chứng AI
    \item 97/97 kiểm thử đạt, chứng minh tính khả thi
  \end{itemize}
  
  \vspace{1em}
  \centering
  \textbf{Trải nghiệm:} https://mxuanvan02.github.io/ainatomy/
\end{frame}

\begin{frame}{Kiến nghị}
  \begin{enumerate}
    \item \textbf{Bộ GD\&ĐT:} Khuyến khích công cụ AI offline, đặc biệt vùng khó khăn
    \item \textbf{Trường THPT:} Tích hợp AInatomy vào Tin học 10, tập huấn giáo viên
    \item \textbf{Giáo viên:} Sử dụng AInatomy, khuyến khích học sinh质疑 AI
    \item \textbf{Nghiên cứu tiếp:} Đánh giá tác động dài hạn, mở rộng khối 11--12
  \end{enumerate}
\end{frame}

\begin{frame}{Tài liệu tham khảo}
  \small
  \begin{enumerate}
    \item Bộ GD\&ĐT (2018). \textit{Chương trình GDPT -- Môn Tin học}.
    \item Bộ GD\&ĐT (2024). \textit{QĐ 2422/QĐ-BGDĐT về giáo dục AI}.
    \item Nguyen \& Le (2023). Teaching AI literacy in high schools. \textit{J. Ed. Tech.}, 15(2), 45--62.
    \item Smith \& Johnson (2024). AI verification skills for students. \textit{Int. J. AI Ed.}, 8(1), 112--128.
    \item Tran (2025). AInatomy: Offline AI lab for high school. \textit{Proc. Vietnam Conf. Ed. Tech.}, 78--85.
  \end{enumerate}
\end{frame}

\end{document}
